\section{素理想与极大理想}\label{sec:prime-maximal-ideals}

给定理想 $I\subseteq R$, 商环 $R/I$ 的环论性质怎样反映为理想 $I$ 的性质? 例如, 若 $R/I$ 是整环或域, 这又说明 $I$ 有什么性质? 本节引入两类对环论和代数几何都很重要的特殊理想, 即素理想和极大理想, 借此讨论这些问题. 我们会给出多项式环中素理想与极大理想的若干例子, 最后举例说明怎样用这些概念研究不可约分解.

\begin{Definition}{素理想与极大理想}\label{def:prime-maximal-ideals}
  若理想 $I\subseteq R$ 满足 $I\ne R$, 且对所有 $a,b\in R$ 都有
  \[
    ab\in I\implies a\in I\quad\text{或}\quad b\in I,
  \]
  则称 $I$ 为\textbf{素理想}.

  若理想 $I\subseteq R$ 满足 $I\ne R$, 且不存在理想 $J\subseteq R$ 使得
  \[
    I\subsetneq J\subsetneq R,
  \]
  则称 $I$ 为\textbf{极大理想}.
\end{Definition}

\begin{Example}{素理想名称由来的例子}\label{ex:prime-ideal-eponymous}
  从定义可知 (Exercise~\ref{ex:prime-element-and-principal-ideal}), 非零元素 $p\in R$ 是素元, 当且仅当主理想 $\langle p\rangle\subseteq R$ 是素理想. 特别地, $\langle2\rangle$, $\langle3\rangle$, $\langle5\rangle$ 和 $\langle7\rangle$ 都是 $\mathbb{Z}$ 中的素理想, 但 $\langle1\rangle$, $\langle4\rangle$, $\langle6\rangle$ 和 $\langle9\rangle$ 不是.
\end{Example}

\begin{Example}{{\textmd{$\langle x\rangle\subseteq K[x]$} 是极大理想}}\label{ex:linear-ideal-maximal}
  要证明 $\langle x\rangle$ 极大, 设 $J\subseteq K[x]$ 是理想且 $\langle x\rangle\subsetneq J$; 我们证明 $J=K[x]$.

  由于 $\langle x\rangle$ 由所有不含常数项的多项式组成, $J$ 必至少包含一个常数项非零的多项式:
  \[
    f=a_{0}+a_{1}x+a_{2}x^{2}+\cdots+a_{n}x^{n}\in J
    \quad\text{且}\quad a_{0}\ne0.
  \]
  由于 $\langle x\rangle$ 是 $J$ 的子集, 所有不含常数项的多项式都属于 $J$. 因此
  \[
    g=a_{1}x+a_{2}x^{2}+\cdots+a_{n}x^{n}\in J.
  \]
  因为理想对减法封闭, 得 $a_{0}=f-g\in J$. 又因为 $a_{0}$ 是 $K[x_{1},\ldots,x_{n}]$ 中的可逆元, 所以 $J=K[x_{1},\ldots,x_{n}]$.
\end{Example}

下面的结果用商环的环论性质给出素理想与极大理想的一种有用刻画.

\begin{Proposition}{按素理想和极大理想取商}\label{prop:prime-maximal-quotient-characterization}
  设 $I\subseteq R$ 是理想.
  \begin{enumerate}[label=(\roman*)]
    \item $I$ 是素理想, 当且仅当 $R/I$ 是整环.
    \item $I$ 是极大理想, 当且仅当 $R/I$ 是域.
  \end{enumerate}
\end{Proposition}

\begin{Proof}
  我们证明每一条的正向蕴涵, 逆向蕴涵分别留在 Exercises~\ref{ex:domain-quotient-implies-prime} 和 \ref{ex:field-quotient-implies-maximal}.

  对 (i) 的正向蕴涵, 我们证明其逆否命题. 假设 $R/I$ 不是整环, 于是存在非零陪集 $[a],[b]\in R/I$, 使得 $[a][b]=0$. 根据商环的定义, 这意味着 $a,b\notin I$, 但 $ab\in I$. 换言之, $I$ 不是素理想.

  要证明 (ii) 的正向蕴涵, 设 $I$ 是极大理想, 并令 $[a]\in R/I$ 是一个非零陪集, 即 $a\notin I$. 我们必须证明 $[a]$ 有乘法逆元. 考虑由 $I$ 和 $a$ 生成的理想:
  \[
    J=\langle I,a\rangle.
  \]
  由于 $a\notin I$, 有 $I\subsetneq J$. 由 $I$ 的极大性, 得 $J=R$, 所以 $1\in J$. 根据理想生成集的定义, 可以把 $1$ 写成
  \[
    1=r_{1}b_{1}+\cdots+r_{n}b_{n}+sa
  \]
  其中 $b_{1},\ldots,b_{n}\in I$, 且 $r_{1},\ldots,r_{n},s\in R$. 由于 $I$ 本身是理想, 这意味着 $1=b+sa$, 其中 $b=r_{1}b_{1}+\cdots+r_{n}b_{n}\in I$, 且 $s\in R$. 因此, 在商环 $R/I$ 中,
  \[
    1=[1]=[b+sa]=[b]+[s][a]=[s][a].
  \]
  所以 $[a]$ 有乘法逆元.
\end{Proof}

由于每个域都是整环, 由 Proposition~\ref{prop:prime-maximal-quotient-characterization} 立即得到以下推论.

\begin{Corollary}{极大理想必为素理想}\label{cor:maximal-implies-prime}
  每个极大理想都是素理想.
\end{Corollary}

\begin{Example}{{\textmd{$\langle x^{2}\rangle\subseteq R[x]$} 不是素理想}}\label{ex:square-variable-ideal-not-prime}
  注意,
  \[
    [x]\in\frac{R[x]}{\langle x^{2}\rangle}
  \]
  是商环中的非零元素, 但 $[x][x]=[x^{2}]=0$. 因此 $[x]$ 是零因子. 由于商环 $R[x]/\langle x^{2}\rangle$ 有零因子, 它不是整环. 由 Proposition~\ref{prop:prime-maximal-quotient-characterization} 可知, $\langle x^{2}\rangle$ 不是素理想. 特别地, 这意味着 $\langle x^{2}\rangle$ 不是极大理想, 也可直接验证:
  \[
    \langle x^{2}\rangle\subsetneq\langle x\rangle\subsetneq K[x].
  \]
\end{Example}

\begin{Example}{{\textmd{$\langle y-x^{2}\rangle\subseteq K[x,y]$} 是素理想但不是极大理想}}
  \leavevmode\label{ex:parabola-ideal-prime-not-maximal}\strut\par\nopagebreak[4]\noindent
  由 Example~\ref{ex:quotient-parabola-isomorphism},
  \[
    K[x,y]/\langle y-x^{2}\rangle\cong K[x].
  \]
  由于 $K[x]$ 是整环但不是域, 可知 $\langle y-x^{2}\rangle$ 是素理想, 但不是极大理想.
\end{Example}

\begin{Example}{{\textmd{$\langle x_{1},\ldots,x_{n}\rangle\subseteq K[x_{1},\ldots,x_{n}]$} 是极大理想}}
  \leavevmode\label{ex:variables-ideal-maximal}\strut\par\nopagebreak[4]\noindent
  由 Example~\ref{ex:quotient-evaluation-origin},
  \[
    K[x_{1},\ldots,x_{n}]/\langle x_{1},\ldots,x_{n}\rangle\cong K.
  \]
  由于 $K$ 是域, 可知 $\langle x_{1},\ldots,x_{n}\rangle\subseteq K[x_{1},\ldots,x_{n}]$ 是极大理想.
\end{Example}

\begin{Example}{{\textmd{$\langle x^{2}+1\rangle\subseteq\mathbb{R}[x]$} 与 \textmd{$\langle x^{2}+1\rangle\subseteq\mathbb{C}[x]$}}}
  \leavevmode\label{ex:quadratic-ideal-real-vs-complex}\strut\par\nopagebreak[4]\noindent
  考虑商环
  \[
    \frac{\mathbb{R}[x]}{\langle x^{2}+1\rangle}.
  \]
  注意, $[x]$ 满足 $[x]^{2}=-[1]=-1$. 换言之, $[x]$ 是 $-1$ 的一个平方根. 我们还知道另一个有 $-1$ 的平方根的环, 即复数域 $\mathbb{C}$.

  考虑函数
  \[
    \begin{aligned}
      \phi:\mathbb{R}[x]&\to\mathbb{C}\\
      f(x)&\mapsto f(\mathrm{i}).
    \end{aligned}
  \]
  可以证明 (Exercise~\ref{ex:real-quadratic-quotient-complex}), $\phi$ 是满的环同态, 且 $\ker(\phi)=\langle x^{2}+1\rangle$. 因此, 由第一同构定理, 可知这个商环与复数域同构:
  \[
    \frac{\mathbb{R}[x]}{\langle x^{2}+1\rangle}\cong\mathbb{C}.
  \]
  所以 $\langle x^{2}+1\rangle\subseteq\mathbb{R}[x]$ 是极大理想.

  现在考虑商环
  \[
    \frac{\mathbb{C}[x]}{\langle x^{2}+1\rangle}.
  \]
  在这个情形下, $[x-\mathrm{i}]$ 和 $[x+\mathrm{i}]$ 都不是零, 但其乘积是零:
  \[
    [x-\mathrm{i}][x+\mathrm{i}]=[x^{2}+1]=0\in\frac{\mathbb{C}[x]}{\langle x^{2}+1\rangle}.
  \]
  因此, $\mathbb{C}$ 上的商环有零因子, 不是整环. 所以 $\langle x^{2}+1\rangle\subseteq\mathbb{C}[x]$ 既不是素理想也不是极大理想.
\end{Example}

最后给出一个应用, 说明素理想和极大理想的概念怎样帮助我们研究不可约元与素元. 先定义一类性质特别好的环.

\begin{Definition}{主理想整环}\label{def:principal-ideal-domain}
  如果整环 $R$ 中的每个理想都是主理想, 则称 $R$ 为\textbf{主理想整环} (PID).
\end{Definition}

例如, 整数环 $\mathbb{Z}$ 是 PID, 读者不妨在 Exercise~\ref{ex:integers-pid} 中验证. 如 Exercise~\ref{ex:variables-ideal-not-principal} 所示, 当 $n>1$ 时, 多项式环 $K[x_{1},\ldots,x_{n}]$ 不是 PID, 因为 $\langle x_{1},\ldots,x_{n}\rangle$ 不是主理想. 不过, 正如下一节将看到的, 一元多项式环 $K[x]$ 是 PID.

下面的结果利用素理想和极大理想的概念证明每个 PID 都是 UFD. 特别地, 这一结果与 Exercise~\ref{ex:integers-pid} 一起, 给出了 $\mathbb{Z}$ 中素因数分解唯一性的自足证明.

\begin{Proposition}{PID 都是 UFD}\label{prop:pid-is-ufd}
  每个主理想整环都是唯一分解整环.
\end{Proposition}

\begin{Proof}
  设 $R$ 是 PID. 先证明 $R$ 是 FD. 为作反证, 假设存在非零且不可逆的 $a\in R$, 不能分解成不可约元的乘积. 这意味着 $a$ 不是不可约元, 因此可将它分解成 $a=a_{1}b_{1}$, 其中 $a_{1}$ 和 $b_{1}$ 都不是可逆元. 如果 $a_{1}$ 和 $b_{1}$ 都能分解成不可约元的乘积, 则 $a$ 也能如此分解. 因而, 不失一般性, 假设 $a_{1}$ 不能分解成不可约元的乘积, 写成 $a_{1}=a_{2}b_{2}$, 其中 $a_{2}$ 和 $b_{2}$ 都不是可逆元. 同前, 不失一般性, 可以假设 $a_{2}$ 不能分解成不可约元的乘积.

  继续上述过程, 递归构造一个序列
  \[
    (a=a_{0},a_{1},a_{2},a_{3},\ldots)
  \]
  其中对每个 $i\ge0$, 都有 $a_{i}=a_{i+1}b_{i+1}$, 且 $b_{i+1}$ 是非可逆元. 由 Exercise~\ref{ex:factor-strict-principal-inclusion}, 可知理想 $\langle a_{i}\rangle$ 构成一个严格包含的链:
  \[
    \langle a_{0}\rangle\subsetneq\langle a_{1}\rangle\subsetneq\langle a_{2}\rangle\subsetneq\langle a_{3}\rangle\subsetneq\cdots.
  \]
  由 Exercise~\ref{ex:ascending-ideal-union}, 上述理想的并是理想. 由于 $R$ 是 PID, 可选 $c\in R$ 使得
  \[
    \langle c\rangle=\bigcup_{i=0}^{\infty}\langle a_{i}\rangle.
  \]
  则必存在某个 $n$ 使 $c\in\langle a_{n}\rangle$, 从而对所有 $k\ge n$ 都有 $\langle c\rangle=\langle a_{k}\rangle$. 这与严格包含 $\langle a_{n}\rangle\subsetneq\langle a_{n+1}\rangle$ 矛盾.

  现在要证明 $R$ 是 UFD. 由 Proposition~\ref{prop:ufd-iff-irreducibles-prime}, 只须证明 $R$ 的每个不可约元都是素元. 设 $p\in R$ 不可约. 要证明 $p$ 是素元, 只须证明 $\langle p\rangle$ 极大. 为此, 设 $I\subseteq R$ 是理想, 且 $\langle p\rangle\subseteq I\subseteq R$. 我们必须证明 $I=\langle p\rangle$ 或 $I=R$.

  由于 $R$ 是 PID, 对某个 $r\in R$ 有 $I=\langle r\rangle$. 因为 $p\in\langle p\rangle\subseteq\langle r\rangle$, 所以对某个 $s\in R$ 有 $p=rs$. 由 $p$ 的不可约性, $r$ 或 $s$ 是可逆元. 若 $r$ 是可逆元, 则 $I=R$; 若 $s$ 是可逆元, 则 $I=\langle p\rangle$.
\end{Proof}

\BookExerciseHeading
\begin{BookExercises}
  \item\label{ex:prime-element-and-principal-ideal} 证明: 非零元素 $p\in R$ 是素元, 当且仅当主理想 $\langle p\rangle\subseteq R$ 是素理想.
  \item\label{ex:two-prime-generators-not-prime} 证明: 即使两个生成元都是素元, $\langle y,y-x^{2}\rangle\subseteq K[x,y]$ 也不是素理想.
  \item\label{ex:zero-ideal-prime-iff-domain} 证明: 零理想 $\langle0\rangle\subseteq R$ 是素理想, 当且仅当 $R$ 是整环.
  \item\label{ex:domain-quotient-implies-prime} 设 $I\subseteq R$ 是理想, 且 $R/I$ 是整环. 证明 $I$ 是素理想.
  \item\label{ex:field-quotient-implies-maximal} 设 $I\subseteq R$ 是理想, 且 $R/I$ 是域. 证明 $I$ 是极大理想.
  \item\label{ex:real-quadratic-quotient-complex} 证明
    \[
      \frac{\mathbb{R}[x]}{\langle x^{2}+1\rangle}\cong\mathbb{C}.
    \]
  \item\label{ex:integers-pid} 证明 $\mathbb{Z}$ 是 PID. (提示: 若 $I\subseteq\mathbb{Z}$ 是理想, 令 $n$ 为 $I$ 中最小的正整数. 证明 $I=\langle n\rangle$.)
  \item\label{ex:factor-strict-principal-inclusion} 设 $R$ 是整环. 若 $a=bc$ 且 $c\notin R^{*}$, 证明 $\langle a\rangle\subsetneq\langle b\rangle$.
  \item\label{ex:ascending-ideal-union} 设 $I_{1}\subseteq I_{2}\subseteq I_{3}\subseteq\cdots$ 是环 $R$ 的理想的一个升链. 证明
    \[
      I=\bigcup_{k=1}^{\infty}I_{k}
    \]
    是 $R$ 的理想.
  \item\label{ex:nonzero-primes-maximal-in-pid} 设 $R$ 是 PID. 证明 $R$ 中每个非零素理想都极大.
\end{BookExercises}

\noindent\textbf{开放获取}\quad 本章按 \href{http://creativecommons.org/licenses/by-nc/4.0/}{Creative Commons Attribution-NonCommercial 4.0 International License} 的条款授权. 只要适当注明原作者与来源, 提供该许可的链接, 并说明是否作过改动, 即可为非商业目的以任何媒介或格式使用, 分享, 改编, 分发和复制.

除非材料的署名说明另有注明, 本章的图片及其他第三方材料也纳入本章的 Creative Commons 许可. 如果某材料未纳入本章的 Creative Commons 许可, 且预期用途不受法律法规许可或超出许可使用范围, 则须直接向版权所有者取得许可.
